% Logros del laboratorio 8 frente a sus objetivos y trabajo del docente.
% Generado por ../generar_logros.py desde ../logros_datos.py: editar los datos y regenerar.
\documentclass[a4paper,11pt]{article}
% Preámbulo común de los informes de laboratorio (II/2026). Se incluye desde laboratorios/labNN_*/
%   % Preámbulo común de los informes de laboratorio (II/2026). Se incluye desde laboratorios/labNN_*/
%   % Preámbulo común de los informes de laboratorio (II/2026). Se incluye desde laboratorios/labNN_*/
%   \input{../informe_comun}
\input{../../gemelo-digital/manual/estilo_manual}
\input{../../gemelo-digital/manual/tecnico/macros}
\graphicspath{{../../documentacion/planos_isa/pid/}{../../documentacion/planos_isa/electrico/}{../../documentacion/manuales_inversa/}
  {../../database/mockups/}{../../database/documentacion/}{../../database/modelo/}{../../database/interfaz/}
  {../../gemelo-digital/manual/img/}{../../modelo-3d/informe/figuras/}{../../modelo-3d/levantamiento/}{../../../../global/fotos-planta/}{../figuras/}}
\usepackage{pdflscape}
\lstset{showstringspaces=false}
% plano A3 en una página apaisada: \plano[página]{archivo}{leyenda}
\newcommand{\plano}[3][1]{\begin{landscape}\begin{figure}[p]\centering\includegraphics[page=#1, width=\linewidth, height=0.86\textwidth, keepaspectratio]{#2}\caption{#3}\end{figure}\end{landscape}}
\setlength{\headheight}{14pt}
\newcommand{\nlab}{}
\newcommand{\tlab}{}
\lhead{\small Brazo XK-JS1-1 · Informe del laboratorio \nlab}
\rhead{\small UMSA · II/2026}

% \portada{número}{título}{inicio}{entrega}{semanas}
\newcommand{\portada}[5]{%
\renewcommand{\nlab}{#1}%
\begin{titlepage}
\centering
\vspace*{0.6cm}
{\large Universidad Mayor de San Andrés · Facultad de Ingeniería\par}
{\large Carrera de Ingeniería Electrónica\par}
\vspace{0.8cm}
{\color{naranja!80!black}\rule{\linewidth}{1.5pt}\par}
\vspace{0.3cm}
{\Large Informe del Laboratorio Integrado Nº #1\par}
\vspace{0.2cm}
{\huge\bfseries #2\par}
\vspace{0.3cm}
{\color{naranja!80!black}\rule{\linewidth}{1.5pt}\par}
\vspace{0.8cm}
{\large Planta: brazo robótico XK-JS1-1 de 4 grados de libertad\par}
\vspace{0.3cm}
{\large ETN-902 Sistemas de Control II · ETN-1034 Aplicación de Técnicas de Control · ETN-1000 Bases de Datos\par}
\vspace{1cm}
\begin{tabular}{rl}
\textbf{Docente:} & Ing. Jorge Antonio Nava Amador\\
\textbf{Auxiliares:} & Limber Ajnota · Fabián Plata · Franz Choque\\
\textbf{Inicio de trabajo:} & #3\\
\textbf{Entrega:} & #4 (#5)\\
\end{tabular}
\vspace{1cm}

\begin{tabular}{ll}
\multicolumn{2}{l}{\textbf{Escuadra del brazo de 4 GDL}}\\[2pt]
Control (ETN-902 / ETN-1034): & Gabriel Andres Gomez Calla\\
 & Jose Lenin Mamani Choque\\
 & Miguel Alejandro Lima Huarachi\\
 & Román Abner Machicado Arequipa\\[2pt]
Cómputo (ETN-1000): & Miguel Alejandro Lima Huarachi\\
 & Marco Eduardo Chura Rondo\\
\end{tabular}
\vfill
{\small Repositorio: \url{https://github.com/GabAndGom/Brazo-de-4-grados-de-libertad}\par}
{\small La Paz, Bolivia · gestión II/2026\par}
\end{titlepage}
}

% Tabla de entregables (Definition of Done): \begin{dod} ... \end{dod}, filas: entregable & estado & evidencia\\
\newenvironment{dod}{\begin{center}\small\begin{longtable}{>{\raggedright}p{5.2cm} >{\raggedright}p{1.8cm} >{\raggedright\arraybackslash}p{7.6cm}}
\toprule \textbf{Entregable} & \textbf{Estado} & \textbf{Evidencia / observación}\\ \midrule \endhead}
{\bottomrule\end{longtable}\end{center}}

\input{../../gemelo-digital/manual/estilo_manual}
% Macros del manual técnico: ladder, estados, código y rutas de imágenes
\usetikzlibrary{shapes.multipart}
\graphicspath{{../img/}{../../../../../global/fotos-planta/}{../../../modelo-3d/informe/figuras/}{../../../modelo-3d/levantamiento/}{../../../database/modelo/}{../../../documentacion/manuales_inversa/}{../../../documentacion/planos_isa/electrico/}{../../../database/interfaz/}}

\lstdefinestyle{py}{language=Python, keywordstyle=\color{azulg}\bfseries, commentstyle=\color{grisos}\itshape,
  stringstyle=\color{verde}, basicstyle=\ttfamily\scriptsize, numbers=left, numberstyle=\tiny\color{grisos}}
\lstdefinestyle{js}{language=Java, morekeywords={const,let,function,async,await,return,for,of,if,else}, keywordstyle=\color{azulg}\bfseries,
  commentstyle=\color{grisos}\itshape, stringstyle=\color{verde}, basicstyle=\ttfamily\scriptsize, numbers=left, numberstyle=\tiny\color{grisos}}
\lstdefinestyle{awl}{basicstyle=\ttfamily\tiny, commentstyle=\color{grisos}\itshape, morecomment=[l]{//},
  keywordstyle=\color{azulg}\bfseries, morekeywords={LD,LDN,A,AN,O,ON,EU,ED,S,R,MOVB,MOVW,MOVD,PLS,TON,CALL,CRET,LPS,LRD,LPP,ALD,OLD,FILL,LDW,LDB,LDD,AW,AD,AB,OW,OD},
  numbers=left, numberstyle=\tiny\color{grisos}, frame=single, backgroundcolor=\color{gris}}
\lstdefinestyle{sql}{language=SQL, keywordstyle=\color{azulg}\bfseries, commentstyle=\color{grisos}\itshape,
  stringstyle=\color{verde}, basicstyle=\ttfamily\tiny, numbers=left, numberstyle=\tiny\color{grisos},
  morekeywords={REFERENCES,SERIAL,SMALLSERIAL,BIGSERIAL,TIMESTAMPTZ,BOOLEAN,INET,RETURNS,TRIGGER,LANGUAGE,DECLARE,RAISE,EXCEPTION,VIEW,INDEX}}

\tikzset{
  ld/.style={line width=0.7pt},
  pics/na/.style={code={\fill[white](-0.12,-0.26) rectangle (0.12,0.26);
     \draw[thick](-0.1,-0.24)--(-0.1,0.24)(0.1,-0.24)--(0.1,0.24);
     \node[font=\scriptsize\ttfamily, above] at (0,0.22){#1};}},
  pics/nc/.style={code={\fill[white](-0.12,-0.26) rectangle (0.12,0.26);
     \draw[thick](-0.1,-0.24)--(-0.1,0.24)(0.1,-0.24)--(0.1,0.24);\draw(-0.17,-0.24)--(0.17,0.24);
     \node[font=\scriptsize\ttfamily, above] at (0,0.22){#1};}},
  pics/pf/.style={code={\fill[white](-0.12,-0.26) rectangle (0.12,0.26);
     \draw[thick](-0.1,-0.24)--(-0.1,0.24)(0.1,-0.24)--(0.1,0.24);
     \node[font=\tiny\bfseries] at (0,0){P};}},
  pics/cmp/.style n args={3}{code={\fill[white](-0.12,-0.26) rectangle (0.12,0.26);
     \draw[thick](-0.1,-0.24)--(-0.1,0.24)(0.1,-0.24)--(0.1,0.24);
     \node[font=\tiny\ttfamily] at (0,0){#2};
     \node[font=\scriptsize\ttfamily, above] at (0,0.22){#1};
     \node[font=\scriptsize\ttfamily, below] at (0,-0.22){#3};}},
  pics/bob/.style n args={3}{code={\fill[white](-0.2,-0.26) rectangle (0.2,0.26);
     \draw[thick](-0.05,0.24) arc (110:250:0.255);\draw[thick](0.05,0.24) arc (70:-70:0.255);
     \node[font=\tiny\bfseries] at (0,0){#2};
     \node[font=\scriptsize\ttfamily, above] at (0,0.22){#1};
     \node[font=\scriptsize\ttfamily, below] at (0,-0.22){#3};}},
  fb/.style={draw, fill=white, font=\scriptsize\ttfamily, align=left, inner sep=3pt, anchor=west},
  barra/.style={line width=2pt},
  estado/.style={draw, rounded corners=3pt, fill=white, align=center, font=\scriptsize, minimum width=22mm, minimum height=10mm, inner sep=3pt},
  estadon/.style={estado, fill=naranja!14, draw=naranja!70!black},
  trans/.style={font=\scriptsize, align=center, fill=white, inner sep=1pt},
  lugar/.style={circle, draw, thick, minimum size=7mm, fill=white},
  transicion/.style={rectangle, fill=black, minimum width=1.6mm, minimum height=7mm, inner sep=0},
  tabla/.style={rectangle split, rectangle split parts=2, draw, fill=white, font=\tiny\ttfamily, align=left, inner sep=2pt, rounded corners=1pt, text width=33mm},
}
\newcommand{\pendiente}[1]{\par\smallskip\noindent\colorbox{naranja!12}{\parbox{\dimexpr\linewidth-2\fboxsep}{\textbf{\color{naranja!80!black}Pendiente.} #1}}\par\smallskip}
\newcommand{\cumple}{\textcolor{verde}{\textbf{cumplido}}}
\newcommand{\parcial}{\textcolor{naranja}{\textbf{parcial}}}
\newcommand{\falta}{\textcolor{rojo}{\textbf{pendiente}}}

\graphicspath{{../../documentacion/planos_isa/pid/}{../../documentacion/planos_isa/electrico/}{../../documentacion/manuales_inversa/}
  {../../database/mockups/}{../../database/documentacion/}{../../database/modelo/}{../../database/interfaz/}
  {../../gemelo-digital/manual/img/}{../../modelo-3d/informe/figuras/}{../../modelo-3d/levantamiento/}{../../../../global/fotos-planta/}{../figuras/}}
\usepackage{pdflscape}
\lstset{showstringspaces=false}
% plano A3 en una página apaisada: \plano[página]{archivo}{leyenda}
\newcommand{\plano}[3][1]{\begin{landscape}\begin{figure}[p]\centering\includegraphics[page=#1, width=\linewidth, height=0.86\textwidth, keepaspectratio]{#2}\caption{#3}\end{figure}\end{landscape}}
\setlength{\headheight}{14pt}
\newcommand{\nlab}{}
\newcommand{\tlab}{}
\lhead{\small Brazo XK-JS1-1 · Informe del laboratorio \nlab}
\rhead{\small UMSA · II/2026}

% \portada{número}{título}{inicio}{entrega}{semanas}
\newcommand{\portada}[5]{%
\renewcommand{\nlab}{#1}%
\begin{titlepage}
\centering
\vspace*{0.6cm}
{\large Universidad Mayor de San Andrés · Facultad de Ingeniería\par}
{\large Carrera de Ingeniería Electrónica\par}
\vspace{0.8cm}
{\color{naranja!80!black}\rule{\linewidth}{1.5pt}\par}
\vspace{0.3cm}
{\Large Informe del Laboratorio Integrado Nº #1\par}
\vspace{0.2cm}
{\huge\bfseries #2\par}
\vspace{0.3cm}
{\color{naranja!80!black}\rule{\linewidth}{1.5pt}\par}
\vspace{0.8cm}
{\large Planta: brazo robótico XK-JS1-1 de 4 grados de libertad\par}
\vspace{0.3cm}
{\large ETN-902 Sistemas de Control II · ETN-1034 Aplicación de Técnicas de Control · ETN-1000 Bases de Datos\par}
\vspace{1cm}
\begin{tabular}{rl}
\textbf{Docente:} & Ing. Jorge Antonio Nava Amador\\
\textbf{Auxiliares:} & Limber Ajnota · Fabián Plata · Franz Choque\\
\textbf{Inicio de trabajo:} & #3\\
\textbf{Entrega:} & #4 (#5)\\
\end{tabular}
\vspace{1cm}

\begin{tabular}{ll}
\multicolumn{2}{l}{\textbf{Escuadra del brazo de 4 GDL}}\\[2pt]
Control (ETN-902 / ETN-1034): & Gabriel Andres Gomez Calla\\
 & Jose Lenin Mamani Choque\\
 & Miguel Alejandro Lima Huarachi\\
 & Román Abner Machicado Arequipa\\[2pt]
Cómputo (ETN-1000): & Miguel Alejandro Lima Huarachi\\
 & Marco Eduardo Chura Rondo\\
\end{tabular}
\vfill
{\small Repositorio: \url{https://github.com/GabAndGom/Brazo-de-4-grados-de-libertad}\par}
{\small La Paz, Bolivia · gestión II/2026\par}
\end{titlepage}
}

% Tabla de entregables (Definition of Done): \begin{dod} ... \end{dod}, filas: entregable & estado & evidencia\\
\newenvironment{dod}{\begin{center}\small\begin{longtable}{>{\raggedright}p{5.2cm} >{\raggedright}p{1.8cm} >{\raggedright\arraybackslash}p{7.6cm}}
\toprule \textbf{Entregable} & \textbf{Estado} & \textbf{Evidencia / observación}\\ \midrule \endhead}
{\bottomrule\end{longtable}\end{center}}

\input{../../gemelo-digital/manual/estilo_manual}
% Macros del manual técnico: ladder, estados, código y rutas de imágenes
\usetikzlibrary{shapes.multipart}
\graphicspath{{../img/}{../../../../../global/fotos-planta/}{../../../modelo-3d/informe/figuras/}{../../../modelo-3d/levantamiento/}{../../../database/modelo/}{../../../documentacion/manuales_inversa/}{../../../documentacion/planos_isa/electrico/}{../../../database/interfaz/}}

\lstdefinestyle{py}{language=Python, keywordstyle=\color{azulg}\bfseries, commentstyle=\color{grisos}\itshape,
  stringstyle=\color{verde}, basicstyle=\ttfamily\scriptsize, numbers=left, numberstyle=\tiny\color{grisos}}
\lstdefinestyle{js}{language=Java, morekeywords={const,let,function,async,await,return,for,of,if,else}, keywordstyle=\color{azulg}\bfseries,
  commentstyle=\color{grisos}\itshape, stringstyle=\color{verde}, basicstyle=\ttfamily\scriptsize, numbers=left, numberstyle=\tiny\color{grisos}}
\lstdefinestyle{awl}{basicstyle=\ttfamily\tiny, commentstyle=\color{grisos}\itshape, morecomment=[l]{//},
  keywordstyle=\color{azulg}\bfseries, morekeywords={LD,LDN,A,AN,O,ON,EU,ED,S,R,MOVB,MOVW,MOVD,PLS,TON,CALL,CRET,LPS,LRD,LPP,ALD,OLD,FILL,LDW,LDB,LDD,AW,AD,AB,OW,OD},
  numbers=left, numberstyle=\tiny\color{grisos}, frame=single, backgroundcolor=\color{gris}}
\lstdefinestyle{sql}{language=SQL, keywordstyle=\color{azulg}\bfseries, commentstyle=\color{grisos}\itshape,
  stringstyle=\color{verde}, basicstyle=\ttfamily\tiny, numbers=left, numberstyle=\tiny\color{grisos},
  morekeywords={REFERENCES,SERIAL,SMALLSERIAL,BIGSERIAL,TIMESTAMPTZ,BOOLEAN,INET,RETURNS,TRIGGER,LANGUAGE,DECLARE,RAISE,EXCEPTION,VIEW,INDEX}}

\tikzset{
  ld/.style={line width=0.7pt},
  pics/na/.style={code={\fill[white](-0.12,-0.26) rectangle (0.12,0.26);
     \draw[thick](-0.1,-0.24)--(-0.1,0.24)(0.1,-0.24)--(0.1,0.24);
     \node[font=\scriptsize\ttfamily, above] at (0,0.22){#1};}},
  pics/nc/.style={code={\fill[white](-0.12,-0.26) rectangle (0.12,0.26);
     \draw[thick](-0.1,-0.24)--(-0.1,0.24)(0.1,-0.24)--(0.1,0.24);\draw(-0.17,-0.24)--(0.17,0.24);
     \node[font=\scriptsize\ttfamily, above] at (0,0.22){#1};}},
  pics/pf/.style={code={\fill[white](-0.12,-0.26) rectangle (0.12,0.26);
     \draw[thick](-0.1,-0.24)--(-0.1,0.24)(0.1,-0.24)--(0.1,0.24);
     \node[font=\tiny\bfseries] at (0,0){P};}},
  pics/cmp/.style n args={3}{code={\fill[white](-0.12,-0.26) rectangle (0.12,0.26);
     \draw[thick](-0.1,-0.24)--(-0.1,0.24)(0.1,-0.24)--(0.1,0.24);
     \node[font=\tiny\ttfamily] at (0,0){#2};
     \node[font=\scriptsize\ttfamily, above] at (0,0.22){#1};
     \node[font=\scriptsize\ttfamily, below] at (0,-0.22){#3};}},
  pics/bob/.style n args={3}{code={\fill[white](-0.2,-0.26) rectangle (0.2,0.26);
     \draw[thick](-0.05,0.24) arc (110:250:0.255);\draw[thick](0.05,0.24) arc (70:-70:0.255);
     \node[font=\tiny\bfseries] at (0,0){#2};
     \node[font=\scriptsize\ttfamily, above] at (0,0.22){#1};
     \node[font=\scriptsize\ttfamily, below] at (0,-0.22){#3};}},
  fb/.style={draw, fill=white, font=\scriptsize\ttfamily, align=left, inner sep=3pt, anchor=west},
  barra/.style={line width=2pt},
  estado/.style={draw, rounded corners=3pt, fill=white, align=center, font=\scriptsize, minimum width=22mm, minimum height=10mm, inner sep=3pt},
  estadon/.style={estado, fill=naranja!14, draw=naranja!70!black},
  trans/.style={font=\scriptsize, align=center, fill=white, inner sep=1pt},
  lugar/.style={circle, draw, thick, minimum size=7mm, fill=white},
  transicion/.style={rectangle, fill=black, minimum width=1.6mm, minimum height=7mm, inner sep=0},
  tabla/.style={rectangle split, rectangle split parts=2, draw, fill=white, font=\tiny\ttfamily, align=left, inner sep=2pt, rounded corners=1pt, text width=33mm},
}
\newcommand{\pendiente}[1]{\par\smallskip\noindent\colorbox{naranja!12}{\parbox{\dimexpr\linewidth-2\fboxsep}{\textbf{\color{naranja!80!black}Pendiente.} #1}}\par\smallskip}
\newcommand{\cumple}{\textcolor{verde}{\textbf{cumplido}}}
\newcommand{\parcial}{\textcolor{naranja}{\textbf{parcial}}}
\newcommand{\falta}{\textcolor{rojo}{\textbf{pendiente}}}

\graphicspath{{../../documentacion/planos_isa/pid/}{../../documentacion/planos_isa/electrico/}{../../documentacion/manuales_inversa/}
  {../../database/mockups/}{../../database/documentacion/}{../../database/modelo/}{../../database/interfaz/}
  {../../gemelo-digital/manual/img/}{../../modelo-3d/informe/figuras/}{../../modelo-3d/levantamiento/}{../../../../global/fotos-planta/}{../figuras/}}
\usepackage{pdflscape}
\lstset{showstringspaces=false}
% plano A3 en una página apaisada: \plano[página]{archivo}{leyenda}
\newcommand{\plano}[3][1]{\begin{landscape}\begin{figure}[p]\centering\includegraphics[page=#1, width=\linewidth, height=0.86\textwidth, keepaspectratio]{#2}\caption{#3}\end{figure}\end{landscape}}
\setlength{\headheight}{14pt}
\newcommand{\nlab}{}
\newcommand{\tlab}{}
\lhead{\small Brazo XK-JS1-1 · Informe del laboratorio \nlab}
\rhead{\small UMSA · II/2026}

% \portada{número}{título}{inicio}{entrega}{semanas}
\newcommand{\portada}[5]{%
\renewcommand{\nlab}{#1}%
\begin{titlepage}
\centering
\vspace*{0.6cm}
{\large Universidad Mayor de San Andrés · Facultad de Ingeniería\par}
{\large Carrera de Ingeniería Electrónica\par}
\vspace{0.8cm}
{\color{naranja!80!black}\rule{\linewidth}{1.5pt}\par}
\vspace{0.3cm}
{\Large Informe del Laboratorio Integrado Nº #1\par}
\vspace{0.2cm}
{\huge\bfseries #2\par}
\vspace{0.3cm}
{\color{naranja!80!black}\rule{\linewidth}{1.5pt}\par}
\vspace{0.8cm}
{\large Planta: brazo robótico XK-JS1-1 de 4 grados de libertad\par}
\vspace{0.3cm}
{\large ETN-902 Sistemas de Control II · ETN-1034 Aplicación de Técnicas de Control · ETN-1000 Bases de Datos\par}
\vspace{1cm}
\begin{tabular}{rl}
\textbf{Docente:} & Ing. Jorge Antonio Nava Amador\\
\textbf{Auxiliares:} & Limber Ajnota · Fabián Plata · Franz Choque\\
\textbf{Inicio de trabajo:} & #3\\
\textbf{Entrega:} & #4 (#5)\\
\end{tabular}
\vspace{1cm}

\begin{tabular}{ll}
\multicolumn{2}{l}{\textbf{Escuadra del brazo de 4 GDL}}\\[2pt]
Control (ETN-902 / ETN-1034): & Gabriel Andres Gomez Calla\\
 & Jose Lenin Mamani Choque\\
 & Miguel Alejandro Lima Huarachi\\
 & Román Abner Machicado Arequipa\\[2pt]
Cómputo (ETN-1000): & Miguel Alejandro Lima Huarachi\\
 & Marco Eduardo Chura Rondo\\
\end{tabular}
\vfill
{\small Repositorio: \url{https://github.com/GabAndGom/Brazo-de-4-grados-de-libertad}\par}
{\small La Paz, Bolivia · gestión II/2026\par}
\end{titlepage}
}

% Tabla de entregables (Definition of Done): \begin{dod} ... \end{dod}, filas: entregable & estado & evidencia\\
\newenvironment{dod}{\begin{center}\small\begin{longtable}{>{\raggedright}p{5.2cm} >{\raggedright}p{1.8cm} >{\raggedright\arraybackslash}p{7.6cm}}
\toprule \textbf{Entregable} & \textbf{Estado} & \textbf{Evidencia / observación}\\ \midrule \endhead}
{\bottomrule\end{longtable}\end{center}}

\renewcommand{\tipodoc}{Logros del Laboratorio Integrado}
\renewcommand{\tipocab}{Logros del laboratorio}
\begin{document}
\portada{8}{Interoperabilidad industrial, API REST del ERP y cierre del proyecto}{martes 8 de diciembre de 2026}{martes 22 de diciembre de 2026}{semanas 18 a 20}
\tableofcontents
\newpage
\nota{\textbf{Estado al 18/11/2026.} Laboratorio por empezar; hay un informe de avance con el plan de integración. Fase 3: implementación, redes industriales y ERP (cierre). Este documento explica, objetivo por objetivo, qué pide la guía, qué logró la escuadra, qué es cada cosa y dónde está la evidencia; después explica el trabajo del docente: qué debía revisar en cada hito y qué tiene preparado la escuadra para esa revisión.}
\section{Resumen}
\begin{center}\small\begin{tabularx}{\linewidth}{c >{\raggedright}p{2.4cm} >{\raggedright}X c c}
\toprule \textbf{Nº} & \textbf{Mención} & \textbf{Objetivo} & \textbf{Estado} & \textbf{\%}\\ \midrule
1 & ETN-902 / ETN-1034 & Interoperabilidad industrial total (Modbus) & \parcial & 10\%\\
2 & ETN-902 / ETN-1034 & Transferencia industrial y cierre del proyecto & \parcial & 50\%\\
3 & ETN-1000 & Despliegue de la API REST del ERP & \parcial & 25\%\\
4 & ETN-1000 & Consolidación de métricas en línea & \parcial & 30\%\\
\bottomrule\end{tabularx}\end{center}
El porcentaje mide cuánto del objetivo está hecho y comprobado en la planta o en el repositorio; no cuenta lo que solo está planificado.
\section{Logros respecto de los objetivos}
\subsection{Mención Control Automático (ETN-902 / ETN-1034)}
\subsubsection{Interoperabilidad industrial total (Modbus)}
\paragraph{Qué pide la guía.} Unificar los PLC de las 4 plantas en una red Modbus TCP y RTU/RS-485.
\paragraph{Qué se logró.}
\begin{itemize}
\item Comunicación PC $\leftrightarrow$ PLC funcionando por PPI (protocolo nativo del S7-200), documentada byte a byte con su suma de verificación.
\item Plan de paso a Modbus RTU con la librería de STEP 7-Micro/WIN (MBUS\_INIT / MBUS\_SLAVE) y mapa de registros 40001–40010 con los tags de la base.
\end{itemize}
\paragraph{Qué es y por qué importa.} Modbus es un protocolo abierto que hablan casi todos los PLC; con él las cuatro plantas, de distintas marcas, comparten datos en una sola red. RTU va por RS-485 (serie) y TCP por Ethernet.
\paragraph{Evidencia.}
\begin{itemize}\small
\item \codigo{laboratorios/lab08\_modbus\_api\_y\_cierre/informe\_lab08.pdf}
\item \codigo{gemelo-digital/manual/tecnico/manual\_tecnico.pdf}
\end{itemize}
\paragraph{Lo que falta.} Programa Modbus en el PLC, cableado STP al anillo de datos y estadísticas de CRC.
\paragraph{Estado:} \parcial~(10\%).
\subsubsection{Transferencia industrial y cierre del proyecto}
\paragraph{Qué pide la guía.} Validar el sistema bajo operación continua y documentar los manuales finales del SDLC y Kaizen.
\paragraph{Qué se logró.}
\begin{itemize}
\item Manual técnico de 87 páginas, manuales de uso en Linux y Windows, guía completa del brazo, documento de arquitectura y 8 informes de laboratorio.
\item Planos v2.0, programa del PLC, scripts SQL, gemelo y modelo 3D versionados en el repositorio.
\end{itemize}
\paragraph{Qué es y por qué importa.} Es dejar la planta lista para que otro grupo la use y la mantenga: todo documentado, versionado y probado en operación prolongada.
\paragraph{Evidencia.}
\begin{itemize}\small
\item \codigo{gemelo-digital/manual/tecnico/manual\_tecnico.pdf}
\item \codigo{gemelo-digital/manual/}
\end{itemize}
\paragraph{Lo que falta.} Pruebas de estrés en operación continua con cambios de setpoint, y memorias firmadas con ORCID.
\paragraph{Estado:} \parcial~(50\%).
\subsection{Mención Sistemas Computacionales (ETN-1000)}
\subsubsection{Despliegue de la API REST del ERP}
\paragraph{Qué pide la guía.} Servicios web seguros para descargar históricos (GET /api/v1/telemetria/…, /alertas/activas).
\paragraph{Qué se logró.}
\begin{itemize}
\item API de operación funcionando en la app, con JSON y sesión por token: \codigo{GET /api/estado}, \codigo{GET /api/info}, \codigo{POST /api/orden/\{nombre\}}, \codigo{POST /api/mando}.
\item Diseño de la API de históricos con FastAPI (telemetría por planta, alertas activas, ciclos, sesiones) sobre las consultas ya probadas, y del respaldo diario con \codigo{pg\_dump} por cron.
\end{itemize}
\paragraph{Qué es y por qué importa.} Una API REST es un conjunto de direcciones web que devuelven datos en JSON; permite que otro programa (o investigador) descargue los históricos sin entrar a la base. El token impide el acceso anónimo.
\paragraph{Evidencia.}
\begin{itemize}\small
\item \codigo{gemelo-digital/app/}
\item \codigo{laboratorios/lab08\_modbus\_api\_y\_cierre/informe\_lab08.pdf}
\end{itemize}
\paragraph{Lo que falta.} Programar los endpoints de históricos y los respaldos programados.
\paragraph{Estado:} \parcial~(25\%).
\subsubsection{Consolidación de métricas en línea}
\paragraph{Qué pide la guía.} Poner en producción el dashboard con resultados y bitácoras ante la comunidad universitaria.
\paragraph{Qué se logró.}
\begin{itemize}
\item Página pública brazo-4gdl.pages.dev con el modelo 3D del brazo, la documentación y la operación remota por roles; se publica sola con cada cambio del repositorio.
\end{itemize}
\paragraph{Qué es y por qué importa.} Es la vitrina de ciencia abierta: cualquiera puede ver en línea cómo trabaja la planta y sus resultados.
\paragraph{Evidencia.}
\begin{itemize}\small
\item \url{https://brazo-4gdl.pages.dev}
\end{itemize}
\paragraph{Lo que falta.} Dashboard con métricas en vivo leídas de PostgreSQL.
\paragraph{Estado:} \parcial~(30\%).
\section{Trabajo del docente}
La guía asigna al docente el papel de auditor: en cada hito revisa el trabajo de la escuadra y decide si cumple antes de pasar a la etapa siguiente. Para cada tarea se explica qué debía revisar, cuándo, y qué tiene la escuadra para presentarle. El porcentaje es el del Excel \codigo{evaluacion\_tareas\_docente.xlsx}.
\subsection{Evaluación de entrega de proyectos}
\paragraph{Lo que dice la guía.} Presidir el tribunal final y verificar la correspondencia entre los manuales firmados en GitHub y la operación física. \textit{Cuándo:} martes 22 de diciembre de 2026 (semana 20).
\paragraph{Qué debía revisar el docente (hitos de la revisión).}
\begin{itemize}
\item[$\square$] Tomar el manual y comprobar en la planta que cada señal, secuencia y límite descrito se cumple.
\item[$\square$] Lazo cerrado funcionando, red con las otras plantas y telemetría en la base (40 \% de la nota final).
\end{itemize}
\paragraph{Qué tiene la escuadra para esa revisión.}
\begin{itemize}
\item Manual técnico (87 pp.), manuales de uso, planos v2.0, documento de base de datos y 8 informes, verificados con la planta.
\end{itemize}
\paragraph{Lo que falta para que el docente lo apruebe.} Lazo cerrado y red con las otras plantas.
\paragraph{Estado:} \parcial~(70\%).
\subsection{Acreditación del ecosistema ERP}
\paragraph{Lo que dice la guía.} Evaluar el dashboard final como vitrina de ciencia abierta y el cumplimiento del ciclo de vida. \textit{Cuándo:} martes 22 de diciembre de 2026.
\paragraph{Qué debía revisar el docente (hitos de la revisión).}
\begin{itemize}
\item[$\square$] API REST, triggers y dashboard operando como una sola plataforma sin fallas lógicas (30 \%).
\item[$\square$] Trazabilidad del repositorio desde el laboratorio 1 hasta el cierre (30 \%).
\end{itemize}
\paragraph{Qué tiene la escuadra para esa revisión.}
\begin{itemize}
\item Página pública con modelo 3D y operación remota por roles; API de operación en JSON en la app.
\end{itemize}
\paragraph{Lo que falta para que el docente lo apruebe.} API REST de históricos, respaldos y dashboard.
\paragraph{Estado:} \parcial~(20\%).
\subsection{Visitas de egresados (tarea de la fase 3)}
\paragraph{Lo que dice la guía.} Coordinar que egresados del IIDEPROQ/INIAM visiten el laboratorio. \textit{Cuándo:} fase 3, los martes.
\paragraph{Qué debía revisar el docente (hitos de la revisión).}
\begin{itemize}
\item[$\square$] Actividad propia del docente; no se evalúa a la escuadra.
\end{itemize}
\paragraph{Estado:} \textit{no aplica}~(—).
\subsection{Hito de la fase}
\textbf{Hito final (semana 20):} las 4 plantas en lazo cerrado, interconectadas en red industrial, guardando telemetría en tiempo real y mostrándola en un dashboard auditable. Para el brazo falta lo central: realimentación de posición, red Modbus, inserción en la base y dashboard. Lo que ya existe (documentación, planos, PLC, base de datos probada, gemelo y operación remota) es la base sobre la que se construye ese cierre.
\section{Conclusión}
Con avance parcial: Interoperabilidad industrial total (Modbus), Transferencia industrial y cierre del proyecto, Despliegue de la API REST del ERP, Consolidación de métricas en línea.
\par\smallskip Pendientes principales:
\begin{enumerate}
\item Programa Modbus en el PLC, cableado STP al anillo de datos y estadísticas de CRC.
\item Pruebas de estrés en operación continua con cambios de setpoint, y memorias firmadas con ORCID.
\item Programar los endpoints de históricos y los respaldos programados.
\item Dashboard con métricas en vivo leídas de PostgreSQL.
\end{enumerate}
\par\smallskip\noindent Repositorio: \url{https://github.com/GabAndGom/Brazo-de-4-grados-de-libertad}
\end{document}
